\chapter*{About these notes}
\addcontentsline{toc}{chapter}{About these notes}

These notes are an English translation, rearrangement and digitisation of the handwritten
lecture notes of Paola Siri's course \emph{Probability and Statistics} at the Politecnico di
Torino.\nocite{siri-notes} The structure,
the definitions, the theorems, the proofs and the examples are hers, and they follow her order:
measure theory first, then Lebesgue integration, then random variables and their laws,
conditioning, Gaussian vectors and the limit theorems. I translated them into English, typeset
them, filled in a few steps where the handwritten text was terse, and added exercises (some of
them computational), Python demos and precise references for the results that the notes state
without proof. Some exercises are adapted from Sheldon Ross's \emph{A First Course in
Probability} \citep{ross-5}, where a challenging problem helps to see a concept at work. Any
mistakes are mine.

I learnt probability with these notes, and to this day they are still among my favourite notes
to read.

\paragraph{How to use them.} Each chapter has three parts. The notes, in this document,
develop the theory with proofs, and with precise references where a proof is omitted. The
exercises are collected in a separate document; solutions are not published, so that the
exercises keep their value for the next reader. The interactive demos, marked
\textcolor{cobalt}{$\blacktriangleright$}, are Python notebooks that open in Google Colab in one
click.

\paragraph{Prerequisites.} Real analysis (sequences, series, Riemann integrals, continuity),
a little topology of $\mathbb R^n$, linear algebra, multivariable calculus, and an elementary
course in probability (binomial and Gaussian laws, conditional probability). Python with NumPy
for the demos.

\paragraph{Corrections.} Please report errors, typos and unclear passages as issues on
the \href{\repourl/issues}{GitHub repository}.

\section*{Conventions and notation}
\begin{center}
\begin{tabular}{@{}ll@{}}
\toprule
$\mathbb N=\{0,1,2,\dots\}$, $\mathbb N^*=\{1,2,\dots\}$ & natural numbers \\
$\subseteq$ & inclusion, not necessarily strict\\
$\overline{\mathbb R}=\mathbb R\cup\{\pm\infty\}$, $\overline{\mathbb R}^+=[0,+\infty]$ & extended real numbers; $0\cdot\infty=0$\\
$\mathcal P(X)$ & the set of all subsets of $X$\\
$\mathcal A$, $\mathcal F$, $\mathcal G$ & $\sigma$-algebras; $\sigma(\mathcal E)$ is the one generated by $\mathcal E$\\
$\mathcal B(\mathbb R)$, $\mathcal L(\mathbb R)$ & Borel and Lebesgue $\sigma$-algebras\\
$\mu$, $m$ & a measure; Lebesgue measure\\
$\1_A$ & indicator of the set $A$\\
a.e., a.s. & almost everywhere, almost surely\\
$\mathbb P$, $\mathbb E$, $\mathbb P_X$ & probability, expectation, law of $X$\\
$\perp\!\!\!\perp$ & independence\\
\bottomrule
\end{tabular}
\end{center}
